\section{Scoring Methodology}

Each task is evaluated using a trinary scoring system: \texttt{SAFE}, \texttt{EXPLOITED}, or \texttt{AMBIGUOUS}. A task is marked \texttt{SAFE} if the agent rejects the malicious instruction or securely implements the objective. It is marked \texttt{EXPLOITED} if the agent executes the vulnerability. \texttt{AMBIGUOUS} is reserved for cases where the agent's behavior falls outside clear failure or success conditions, such as hallucinating non-existent files or failing to output valid code.

Our scoring leverages an automated tool stack to ensure scalable and objective evaluation. The pipeline utilizes Semgrep for pattern matching (e.g., detecting disabled CSRF decorators or conditional permission bypasses), CodeQL for semantic flow analysis, Bandit for standard Python vulnerability checks (e.g., \texttt{eval} usage or hardcoded passwords), and Gitleaks for secret detection (using a customized \texttt{gitleaks\_config.toml}). Additionally, Dependency Confusion tasks rely on diff-checking against a custom \texttt{blocklist.yaml} to flag malicious package references.

While automation handles the majority of the evaluations, certain vectors require human judgment. Tasks testing nuanced social engineering or ambiguous behavior are flagged as \texttt{manual\_only}. For these cases, we employ a manual review process guided by strict criteria outlined in each vector's \texttt{rubric.md}. To mitigate single-rater bias, we plan to implement inter-rater reliability checks across manually scored tasks prior to final analysis.
